\chapter{Maintenance and Development}

\section{Project layout}
\index{project layout}%
\begin{table}[h]
  \centering
  \begin{tabularx}{\textwidth}{lD}
    \toprule
    \textbf{Path} & \textbf{Contents} \\
    \midrule
    \cfg{src/main.rs} & Entry point, UI callbacks, task wiring. \\
    \cfg{src/hardware.rs} & Robot, conveyor and camera drivers plus mocks;
      workspace limits; emergency-stop handles. \\
    \cfg{src/orchestrator.rs} & Message loop, instruction queue, trajectory
      planner. \\
    \cfg{src/vision/} & Detector, calibration, tracker, classifier. \\
    \cfg{src/app\_config.rs} & Configuration structures, loading,
      validation. \\
    \cfg{ui/app\_window.slint} & Operator interface definition. \\
    \cfg{docs/} & Specifications, this manual (sources under
      \cfg{docs/manual/}), TODO list. \\
    \bottomrule
  \end{tabularx}
  \caption{Repository layout}
\end{table}

\section{Tests}
\index{tests}\index{cargo test@\texttt{cargo test}}%
Run the unit test suite with \cfg{cargo test}. Covered today:
configuration parsing and validation (including backward compatibility
with older configuration files), the pixel-to-world calibration transform,
the trajectory-planner intercept mathematics, workspace-limit enforcement
in both robot drivers, and the instruction queue.

\section{Development workflow}
\index{development workflow}\index{mock mode}%
Develop against mock hardware (\cfg{--mock}) whenever possible. The mocks
implement the same traits as the real drivers, including workspace-limit
enforcement, so orchestration logic can be validated safely. The robot's
serial protocol is documented in Appendix~\ref{app:gcode}.

\section{Editing this manual}
\index{manual!editing}\index{latexmk@\texttt{latexmk}}%
The manual sources live under \cfg{docs/manual/}, one file per
chapter; \cfg{docs/manual.tex} is the driver that sets the class
and includes them. Build with:
\begin{lstlisting}[language=bash]
cd docs && latexmk -pdf manual.tex
\end{lstlisting}

\section{Remaining work}
\index{TODO list}\index{roadmap}%
The authoritative, maintained list of open work items lives in
\cfg{docs/TODO.md}. Major open areas at the time of writing:
wiring the vision loop end-to-end (camera $\rightarrow$ detection
$\rightarrow$ pick scheduling), the live camera feed and overlays in the
UI, a real object tracker and classifier, belt-motion compensation using
robot position feedback, and a camera calibration wizard.
